%%%%%%%%%%%%%%%%%%%%%%%%%%%%%%%%%%%%%%%%%%%%%%%%%%%%%%%%%%%%%%%%%%%%%%%%%%%%%%%%
%%  ZWISCHENABGABE: Einlage: Arbeitsplan   (Moodle: 28.10.2026)
%%  Aufbau nach pp-portfolio-geruest-einlage-2026-2027.tex.
%%  Der Text steht in einlagen/1_arbeitsplan.tex und wird 1:1 in die Endabgabe übernommen.
%%  Overleaf: Menü -> Main document = diese Datei. PDF für Moodle benennen:
%%  Kautzsch_Arbeitsplan.pdf
%%%%%%%%%%%%%%%%%%%%%%%%%%%%%%%%%%%%%%%%%%%%%%%%%%%%%%%%%%%%%%%%%%%%%%%%%%%%%%%%
%%%%%%%%%%%%%%%%%%%%%%%%%%%%%%%%%%%%%%%%%%%%%%%%%%%%%%%%%%%%%%%%%%%%%%%%%%%%%%%%
%%  EINSTELLUNGEN – die einzige Datei, in der du deine Angaben einträgst.
%%  Wird von ALLEN Hauptdateien (Endabgabe + Zwischenabgaben) eingelesen.
%%  Entspricht dem Block "Die Angaben" der offiziellen Gerüste
%%  pp-portfolio-geruest-2026-2027.tex / pp-portfolio-geruest-einlage-2026-2027.tex
%%%%%%%%%%%%%%%%%%%%%%%%%%%%%%%%%%%%%%%%%%%%%%%%%%%%%%%%%%%%%%%%%%%%%%%%%%%%%%%%

%%%%%%%%%%%%%%%%%%%%%%%%%%%%%%%%%%%%% Die Angaben %%%%%%%%%%%%%%%%%%%%%%%%%%%%%%%
\newcommand{\Semester}{WH 2026-2027}
\newcommand{\Projekt}{Projekt, Titel des Portfolios}          % z.B. Titel eures Gruppenprojekts
\newcommand{\Schwerpunkt}{Ihr Schwerpunktthema}               % DEIN Schwerpunkt (≠ Partner/in!)
\newcommand{\IhrName}{Moritz Kautzsch}                         % Vorname Nachname
\newcommand{\Matrikelnummer}{Ihre Matrikelnummer}
\newcommand{\Studiengang}{Ihr Studiengang (Name und Stufe)}   % z.B. Physik B.Sc., 5. Semester
\newcommand{\emailAdresse}{Ihre e-Mail-Adresse}               % Uni-Adresse empfohlen
\newcommand{\dasDatum}{Das Datum der Abgabe}                  % z.B. 15.\,03.\,2027

%% Weitere Angaben, die in mehreren Einlagen gebraucht werden
\newcommand{\Partner}{Name Partnerin/Partner}                  % für Abgrenzung der Schwerpunkte
\newcommand{\PartnerSchwerpunkt}{Schwerpunkt Partnerin/Partner}
\newcommand{\Laborbuchlink}{https://link-zum-elektronischen-laborjournal}  % Pflicht in Einlage 3!
\newcommand{\Protokolldatum}{TT.MM.2026}   % Einlage 3: gewählter Tag (≠ Tag der Partnerin/des Partners)
\newcommand{\PaperKey}{VORLAGE-Artikel}        % Einlage 2: bib-Key deines Papers in literatur.bib
\newcommand{\Vortragsdatum}{28.\,01.\,2027}    % Einlage 5
\newcommand{\FremdVortragKey}{VORLAGE-Vortrag} % Einlage 6: bib-Key des anderen Vortrags

%%%%%%%%%%%%%%%%%%%%%%%%%%%%%%%%%%%%%%%%%%%%%%%%%%%%%%%%%%%%%%%%%%%%%%%%%%%%%%%%
%%  MODUS-SCHALTER
%%
%%  \Abgabefalse  = Entwurf: grüne Hilfen sichtbar, Kontrollseite mit
%%                  Umfangsmessung, Platzhalter rot markiert.
%%  \Abgabetrue   = Abgabe: grüne Farbe -> schwarz (wie offizielle Zeile 15/16),
%%                  alle Hilfen + Kontrollseite + Umfangsboxen verschwinden.
%%                  Übrig gebliebene \platzhalter{} bleiben ROT sichtbar und
%%                  erzeugen eine Warnung -> dann ist noch etwas offen!
%%
%%  \Hinweisefalse blendet im Entwurf nur die grünen Hilfstexte aus
%%                  (nützlich, um den Umfang realistisch zu messen).
%%%%%%%%%%%%%%%%%%%%%%%%%%%%%%%%%%%%%%%%%%%%%%%%%%%%%%%%%%%%%%%%%%%%%%%%%%%%%%%%
\newif\ifAbgabe
\newif\ifHinweise

\Abgabefalse      % <-- für die Abgabe auf \Abgabetrue setzen
\Hinweisetrue     % <-- Hilfstexte im Entwurf ein-/ausblenden

\newcommand{\dieEinlage}{Einlage: Arbeitsplan}

\documentclass[pdftex,12pt,twoside,footsepline,headsepline,a4paper,final]{scrreprt}
%%%%%%%%%%%%%%%%%%%%%%%%%%%%%%%%%%%%%%%%%%%%%%%%%%%%%%%%%%%%%%%%%%%%%%%%%%%%%%%%
%%  PRÄAMBEL – gemeinsam für Endabgabe und alle Zwischenabgaben.
%%  Teil 1 ist 1:1 die Präambel der offiziellen Gerüste (nicht verändern),
%%  Teil 2 ergänzt nur Pakete für Gleichungen, Bilder, PDF-Einbindung usw.,
%%  Teil 3 enthält die Hilfsmakros dieser Vorlage.
%%  Die \documentclass-Zeile steht (unverändert offiziell) in jeder Hauptdatei.
%%%%%%%%%%%%%%%%%%%%%%%%%%%%%%%%%%%%%%%%%%%%%%%%%%%%%%%%%%%%%%%%%%%%%%%%%%%%%%%%

%%%%%%%%%%%%%%%%%%%%%%%%%%%%% Teil 1: offizielle Präambel %%%%%%%%%%%%%%%%%%%%%%
\usepackage[verbose,a4paper,bottom=25mm,top=25mm]{geometry} %showframe,
\usepackage{scrlayer-scrpage}
\usepackage[T1]{fontenc}
\usepackage[utf8]{inputenc}
\usepackage{ngerman}
\usepackage{color}
\usepackage{amsfonts}
\usepackage{amssymb}
\usepackage{amsthm}
%%%%%%%%%%%%%%%%%%%%%%%%% Biblatex (offizielle Optionen, Backend Biber) %%%%%%%%%%
%% Overleaf erkennt biber automatisch. Zitierstil alphabetic: [Sch02]
\usepackage[sorting=nyt, backend=biber, citestyle=alphabetic, style=alphabetic,
            autocite=inline, labelalpha=true, hyperref=auto, backref=true,
            language=ngerman]{biblatex}
\addbibresource{literatur.bib}

%%%%%%%%%%%%%%%%%%%%%%%%%%%%% Teil 2: Zusatzpakete %%%%%%%%%%%%%%%%%%%%%%%%%%%%%%
\usepackage{amsmath}               % \eqref, align, nummerierte Gleichungen
\usepackage{xcolor}
\definecolor{darkgreen}{rgb}{0,0.5,0}   % offizielle Hinweisfarbe
\usepackage{graphicx}              % \includegraphics
\graphicspath{{abbildungen/}}
\usepackage{booktabs}              % schöne Tabellen
\usepackage[locale=DE]{siunitx}    % \SI{3.4}{\kilo\hertz}, \num{1.2e3}
\usepackage{pdfpages}              % \includepdf für Anhänge (Paper, Laborjournal, Folien)
\usepackage{pdflscape}             % Querformat: \begin{landscape}...\end{landscape}
\usepackage{pgfgantt}              % Gantt-Diagramm für den Arbeitsplan
\usepackage{environ}
\usepackage[savepos,abspage]{zref}
\usepackage[hidelinks]{hyperref}   % klickbare Verweise; zuletzt laden

%% Zitierregel (Marti): bei Journalartikeln DOI angeben, aber keine Webadresse.
%% -> url/urldate werden bei @article automatisch unterdrückt.
\AtEveryBibitem{%
  \ifentrytype{article}{\clearfield{url}\clearfield{urlyear}%
    \clearfield{urlmonth}\clearfield{urlday}}{}}

%% Kopf- und Fußzeile wie offizielles Gerüst (\ohead wird in der Hauptdatei gesetzt)
\pagestyle{scrheadings}
\ihead{\thepage}
\ofoot{\pagemark}
\ifoot{\zubearbeiten{\IhrName}}

%%%%%%%%%%%%%%%%%%%%%%%%%%%%% Teil 3: Hilfsmakros %%%%%%%%%%%%%%%%%%%%%%%%%%%%%%%
\makeatletter

%% --- offizieller Farbschalter (Zeile 12-15 der Gerüste), automatisch per Modus
\ifAbgabe
  \Hinweisefalse
  \newcommand{\zubearbeiten}[1]{\textcolor{black}{#1}}
\else
  \newcommand{\zubearbeiten}[1]{\textcolor{darkgreen}{#1}}
\fi

%% --- Grüne Hilfstexte: verschwinden im Abgabemodus vollständig
%%     \begin{hinweise} ... \end{hinweise}   (darf itemize enthalten)
%%     \hinweis{kurzer Text}
\NewEnviron{hinweise}{\ifHinweise\par{\color{darkgreen}\small\BODY\par}\fi}
\newcommand{\hinweis}[1]{\ifHinweise\par{\color{darkgreen}\small #1\par}\fi}

%% --- Platzhalter für eigenen Text: bleibt IMMER rot sichtbar (auch im
%%     Abgabemodus), damit nichts Unfertiges abgegeben wird.
\newcounter{platzhalter}
\DeclareRobustCommand{\platzhalter}[1]{%
  \stepcounter{platzhalter}%
  \textcolor{red}{\textbf{[}\,#1\,\textbf{]}}%
  \ifAbgabe\PackageWarning{Portfolio}{Platzhalter noch vorhanden: #1}\fi}
\AtEndDocument{%
  \ifnum\value{platzhalter}>0
    \PackageWarningNoLine{Portfolio}{Es sind noch \theplatzhalter\space
      Platzhalter im Dokument!}%
  \fi
  \immediate\write\@auxout{\string\gdef\string\pf@platzhalteralt{\theplatzhalter}}}
\providecommand{\pf@platzhalteralt}{?}

%% --- Umfangsmessung
%%   \begin{umfang}{Bezeichnung}{min}{max} ... \end{umfang}
%%   misst die Länge des eingeschlossenen Textes in Seiten (über die
%%   Position auf der Seite, auf 0.05 Seiten genau) und zeigt sie im Entwurf
%%   in einem Kasten an. Achtung: Abbildungen/Tabellen zählen nach den Regeln
%%   NICHT zur Länge, werden hier aber mitgemessen -> dann abziehen.
\zref@addprop{savepos}{abspage}
\newcounter{pf@umfang}
\newcommand*{\pf@umfangliste}{}
\newcommand*{\pf@umfangeintrag}[4]{% Name, min, max, ist
  \g@addto@macro\pf@umfangliste{\pf@umfangzeile{#1}{#2}{#3}{#4}}}
\newcommand*{\pf@setfarbe}[3]{% min max ist -> \pf@farbe
  \def\pf@farbe{darkgreen}%
  \ifdim #3pt<#1pt \def\pf@farbe{red}\fi
  \ifdim #3pt>#2pt \def\pf@farbe{red}\fi}
\newcommand*{\pf@umfangzeile}[4]{%
  #1 & #2 -- #3 & \pf@setfarbe{#2}{#3}{#4}\textcolor{\pf@farbe}{\textbf{#4}}\\}
\newenvironment{umfang}[3]{%
  \stepcounter{pf@umfang}%
  \edef\pf@uid{pfumf\the\value{pf@umfang}}%
  \def\pf@uname{#1}\def\pf@umin{#2}\def\pf@umax{#3}%
  \par\zsavepos{\pf@uid a}\nobreak
}{%
  \par\zsavepos{\pf@uid e}%
  \ifAbgabe\else
    \zref@refused{\pf@uid a}\zref@refused{\pf@uid e}%
    \edef\pf@ist{\fpeval{round(
        (\zref@extractdefault{\pf@uid e}{abspage}{0}
        -\zref@extractdefault{\pf@uid a}{abspage}{0})
        +(\zref@extractdefault{\pf@uid a}{posy}{0}
        -\zref@extractdefault{\pf@uid e}{posy}{0})/\number\textheight ,2)}}%
    \protected@write\@auxout{}{\string\pf@umfangeintrag{\pf@uname}{\pf@umin}{\pf@umax}{\pf@ist}}%
    \par\smallskip\noindent
    \pf@setfarbe{\pf@umin}{\pf@umax}{\pf@ist}%
    \fcolorbox{\pf@farbe}{white}{\scriptsize
      Umfang \emph{\pf@uname}: $\approx$\,\pf@ist~Seiten
      (Soll: \pf@umin\,--\,\pf@umax~Seiten, ohne Titel/Abb./Tab./Legenden)}%
    \par\smallskip
  \fi}

%% --- Kontrollseite (nur im Entwurf; ersetzt das grüne Kapitel
%%     "Hinweise zur Arbeit mit dieser Datei", das abgegeben nicht vorkommen darf)
\newcommand{\kontrollseite}{%
  \ifAbgabe\else
  \chapter*{Kontrollseite (erscheint nur im Entwurf)}
  {\color{darkgreen}
  Modus: \textbf{Entwurf}. Für die Abgabe in \texttt{einstellungen.tex}
  \texttt{\textbackslash Abgabetrue} setzen. Dann wird alles Grüne schwarz bzw.\
  entfernt (Regel: \enquote{Abgaben mit grüner Farbe oder stehen gebliebenen
  Hilfen ergeben einen Abzug}).\par\medskip
  \textbf{Offene Platzhalter} (letzter Durchlauf): \pf@platzhalteralt\par\medskip
  \textbf{Umfangskontrolle} (Werte aus dem letzten Kompilierdurchlauf;
  Hilfstexte stehen außerhalb der Messung):\par\smallskip
  \ifx\pf@umfangliste\@empty
    \emph{noch keine Messwerte -- nochmals kompilieren.}
  \else
    {\small\begin{tabular}{p{0.55\textwidth}cc}
      \toprule Teil & Soll [Seiten] & Ist [Seiten]\\\midrule
      \pf@umfangliste
      \bottomrule
    \end{tabular}}
  \fi\par\medskip
  %% Checkliste – erscheint nur auf der Kontrollseite im Entwurf.
%% Quellen: Portfolios im Projektpraktikum 2026-2027 (Marti, 09.07.2026),
%% offizielle Gerüste 2026-2027, Zitierregeln (31.08.2026),
%% Abbildungen und Gleichungen (Limmer 2007), Moodle-Terminliste.

\textbf{Fristen (Moodle; nächste Abgabe wird erst nach der vorherigen freigeschaltet!)}\par\smallskip
{\small\begin{tabular}{@{}ll@{}}
21.\,10.\,2026 & Motivationsformular (ohne Reflexion) -- obligatorisch\\
28.\,10.\,2026 & Arbeitsplan + Reflexion -- obligatorisch\\
11.\,11.\,2026 & Wissenschaftliche Arbeit + Reflexion -- obligatorisch\\
09.\,12.\,2026 & Protokoll eines Praktikumstages + Reflexion -- obligatorisch\\
20.\,01.\,2027 & Messergebnis + Reflexion -- obligatorisch\\
23.\,01.\,2027 & Späteste Anmeldung mündliche Leistung (Portal)\\
27.\,01.\,2027 & Folien des Vortrags + Materialien für Mitstudierende\\
28.\,01.\,2027 & Vorträge (Prüfung, Anwesenheitspflicht)\\
07.\,02.\,2027 & Folien mit Kommentar + Reflexion (vorgeschlagen, freiwillig)\\
21.\,02.\,2027 & Aussage aus anderem Vortrag + Reflexion (vorgeschlagen, freiwillig)\\
10.\,03.\,2027 & Späteste Anmeldung schriftliche Leistung (Portal)\\
\textbf{15.\,03.\,2027} & \textbf{Endabgabe Portfolio -- verspätet = nicht bestanden}\\
\end{tabular}}\par\medskip

\textbf{Vor jeder Abgabe prüfen}
\begin{itemize}\itemsep0pt
\item \texttt{\textbackslash Abgabetrue} gesetzt, keine grüne Farbe, keine Hilfen, \emph{0 Platzhalter}.
\item Deckblatt vollständig: Name, Vorname, Matrikelnummer, Studiengang, E-Mail, Thema (Projekt + Schwerpunkt).
\item Inhaltsverzeichnis mit Seitenzahlen, auch der Anhänge; durchgehend Seitenzahlen.
\item Längen eingehalten (siehe Tabelle oben; Abbildungen/Tabellen/Legenden zählen nicht mit):
      Einleitung 2--4, Beschreibung $\tfrac14$--$\tfrac12$, Diskussion der Physik 1--3,
      Reflexion 1--2, Gesamtreflexion 3--6 Seiten. Schriftgröße 12\,pt (fest eingestellt).
\item Beschreibung und Reflexion sind \emph{keine Nacherzählung}; die Leitfragen werden in Fließtext
      beantwortet (nicht in Einzelsätzen), erkennbar zugeordnet.
\item Paper (Einlage 2) wird in den späteren Einlagen zur Argumentation verwendet; Gleichungen werden
      zur Argumentation genutzt.
\item Peer-Feedback (Kommilitonen) eingeholt, aufgegriffen -- und zitiert.
\item Jede Übernahme zitiert: Literatur, Bilder (\texttt{\textbackslash cite[Abb.~3.4]\{key\}}),
      Gleichungen (\texttt{\textbackslash cite[Gl.~(3)]\{key\}}), Aussagen von Personen, Anleitungen,
      Datenblätter, Webseiten (mit Abrufdatum), \textbf{jede KI-Nutzung mit System und ALLEN Prompts}.
\item Journalartikel: Autoren, Titel, Journal, Band, Seiten, Jahr, DOI (nicht als Link).
\item Abbildungen: Bildunterschrift, im Text referenziert (\texttt{\textbackslash ref}), Schrift
      $\geq 10$\,pt, Achsen mit Einheit als \enquote{$f$ / kHz} (nicht \enquote{[kHz]}), ggf.\ Residuenplot.
\item Tabellen: Überschrift, im Text referenziert; Messdaten nicht als lange Listen im Text, sondern im Anhang.
\item Wichtige Gleichungen nummeriert, alle Symbole definiert und einheitlich verwendet.
\item Dateiname für Moodle: \texttt{Kautzsch\_WelcheEinlage.pdf} (nur PDF oder DOCX).
\item Bis auf den Arbeitsplan müssen sich alle Einlagen von denen der Partnerin/des Partners unterscheiden.
\end{itemize}

  }%
  \fi}
\providecommand{\enquote}[1]{\glqq #1\grqq}

%% --- Anhang aus PDF mit durchgehenden Seitenzahlen (Regel: "Die Arbeit ist
%%     durchgängig mit Seitenzahlen versehen ... auch der Anhänge")
%%   \anhangpdf{Kapiteltitel}{label}{datei.pdf}           alle Seiten
%%   \anhangpdf[pages={1-3}]{Kapiteltitel}{label}{datei.pdf}
%%   Die Kapitelüberschrift steht auf einer eigenen Seite, danach folgen die
%%   PDF-Seiten (auf 82 % verkleinert, damit Kopf-/Fußzeile frei bleiben).
\newcommand{\anhangpdf}[4][pages=-]{%
  \chapter{#2}\label{#3}%
  \includepdf[#1,scale=0.82,pagecommand={\thispagestyle{scrheadings}}]{#4}}
%% Gleiches ohne neues Kapitel (für weitere PDFs im selben Anhang)
\newcommand{\pdfseiten}[2][pages=-]{%
  \includepdf[#1,scale=0.82,pagecommand={\thispagestyle{scrheadings}}]{#2}}

\makeatother

%% Antwortfeld im Motivationsformular (Anhang der Endabgabe)
\newcommand{\antwort}[1]{\par\smallskip #1\par\bigskip}

\ohead{\zubearbeiten{\Projekt{} - \Schwerpunkt}}

\begin{document}

%% Titelseite einer ZWISCHENABGABE – wörtlich aus pp-portfolio-geruest-einlage-2026-2027.tex
%% \dieEinlage wird in der jeweiligen Hauptdatei (Kautzsch_....tex) gesetzt.
\title{Universit\"{a}t Ulm\\\vspace{1cm}
{\LARGE \textnormal{\textsc{Fakult\"{a}t f\"{u}r Naturwissenschaften}}}\\\vspace{2cm} {\LARGE \textnormal{\textsl{Projektpraktikum \Semester}}}\\\vspace{2cm} {\LARGE\zubearbeiten{\textit{\Projekt}}}\\
{\Large \zubearbeiten{\textit{\Schwerpunkt}}}\\\vspace{0.5cm} \zubearbeiten{\textit{\dieEinlage}}}

\author{{\LARGE \zubearbeiten{\IhrName}}\\ $_{}$\\$_{}$\\\vspace*{1.cm}
Matrikelnummer:  \zubearbeiten{\Matrikelnummer}\\\vspace{0.8cm} Studiengang: \zubearbeiten{\Studiengang}\\\vspace{0.4cm} e-Mail: \zubearbeiten{\emailAdresse}}

\date{\zubearbeiten{\dasDatum}}

\maketitle

\tableofcontents
\kontrollseite   % nur im Entwurf sichtbar

\chapter{Hauptteil}
\begin{hinweise}
Beachten Sie die Hinweise in der Datei mit dem Gerüst für die Endabgabe! Auch bei einer Zwischenabgabe ist
eine korrekte Literaturarbeit notwendig. Denken Sie daran, dass Sie alle Hinweise im Text auf andere Werke
(Artikel, Bücher, Bilder, Grafiken) aber auch alle Erwähnungen von Meinungen anderer Personen mit einem Zitat
versehen müssen! Beachten Sie auch die Fragen bei der jeweiligen Abgabe in Moodle.
\end{hinweise}

%%%%%%%%%%%%%%%%%%%%%%%%%%%%%%%%%%%%%%%%%%%%%%%%%%%%%%%%%%%%%%%%%%%%%%%%%%%%%%%%
%%  EINLAGE 1: ARBEITSPLAN          Zwischenabgabe: 28.10.2026 (obligatorisch)
%%  Struktur: Beschreibung (1/4–1/2 S.) + Reflexion (1–2 S.)
%%  Die Einlage selbst (der Plan) steht im Anhang: 1_arbeitsplan_anhang.tex
%%%%%%%%%%%%%%%%%%%%%%%%%%%%%%%%%%%%%%%%%%%%%%%%%%%%%%%%%%%%%%%%%%%%%%%%%%%%%%%%
\section{Einlage: Arbeitsplan}
\label{sec:arbeitsplan}

\begin{hinweise}
Diese Einlage ist Ihr Arbeitsplan. Der Arbeitsplan ist die \emph{einzige} Einlage, die mit der Partnerin/dem
Partner identisch sein darf (\enquote{Ihr gemeinsamer Arbeitsplan}). Hinweise zur Erstellung (Marti): nicht nur
eine grobe Liste, sondern feine Teilaufgaben mit zeitlichem Ablauf, sichtbar machen, was parallel läuft;
Vorbereitung (\enquote{Einkaufen}) und Nachbereitung (\enquote{Abwasch}) nicht vergessen -- z.\,B. als
Excel-Tabelle oder in ProjectLibre, oder mit dem Gantt-Gerüst im Anhang.
\end{hinweise}

\subsection{Beschreibung}
\begin{hinweise}
Die Beschreibung gibt alle notwendigen Informationen an, die die Einlage beschreiben, also zum Beispiel den
Fundort, die Literaturstelle und eine kurze Inhaltsübersicht, aber keine Nacherzählung ($1/4$ Seite, max.
$1/2$ Seite!). Leitfaden: Um welche Art von Einlage handelt es sich? Woher stammt sie? Welches sind die
wesentlichen Aussagen?
\end{hinweise}
\begin{umfang}{Einlage 1 (Arbeitsplan): Beschreibung}{0.25}{0.5}
\platzhalter{Art der Einlage (z.\,B. Gantt-Diagramm, erstellt mit \dots), wann/wie gemeinsam mit \Partner{}
erstellt, Zeitraum, Hauptphasen, Meilensteine; Verweis auf Anhang~\ref{anh:arbeitsplan}}
\end{umfang}

\subsection{Reflexion}
\begin{hinweise}
In der Reflexion sollen Sie sich Gedanken machen über die Wirkung dieser Einlage und der Arbeit mit ihr auf
Ihren Lernfortschritt.
\begin{itemize}\itemsep0pt
\item In welchem Zusammenhang steht der Arbeitsplan zu Ihrem Projektthema?
\item In welchem Zusammenhang steht der Arbeitsplan zu Ihrem Schwerpunktthema?
\item Was haben Sie bei der Erstellung des Arbeitsplanes gelernt?
\item Wie hat Ihnen der Arbeitsplan beim Verständnis Ihres Projektthemas und Ihres Schwerpunktthemas geholfen?
\item Wo gab es Schwierigkeiten bei der Beschäftigung mit der Arbeitsplan?
\item Was sagen Ihre Kommilitoninnen oder Kommilitonen dazu, als Kommentar oder wenn Sie mit Ihnen über den Arbeitsplan diskutieren?
\item Enthält Ihr Arbeitsplan logische oder physikalische Fehler?
\item Gab es bei der Beschäftigung mit dem Arbeitsplan eine überraschende neue Erkenntnis?
\end{itemize}
\textbf{Die Reflexion für eine Einlage sollte mindestens eine Seite und maximal 2 Seiten lang sein. Die
Längenangaben sind ohne Titel, Bilder, Tabellen und Tabellen- oder Bildlegenden zu verstehen.}
Peer-Feedback zitieren (siehe \texttt{literatur.bib}, Vorlage \texttt{VORLAGE-Gespraech}).
\end{hinweise}
\begin{umfang}{Einlage 1 (Arbeitsplan): Reflexion}{1}{2}
\platzhalter{Reflexion in Fließtext -- Leitfragen erkennbar beantworten}
\end{umfang}

%% Optional (vom Leitfaden ausdrücklich erwünscht): eigene Bewertungskriterien
%\subsection{Eigene Bewertungskriterien}
%Für diese Einlage schlage ich zusätzlich folgende Kriterien vor: ...


\printbibliography[heading=bibintoc,title={Referenzen und Literaturverzeichnis}]

\appendix
\chapter*{Anh\"{a}nge}
\addcontentsline{toc}{chapter}{Anh\"{a}nge}
%%%%%%%%%%%%%%%%%%%%%%%%%%%%%%%%%%%%%%%%%%%%%%%%%%%%%%%%%%%%%%%%%%%%%%%%%%%%%%%%
%%  ANHANG zu Einlage 1: der Arbeitsplan selbst
%%  Variante A: Plan als PDF (Excel/ProjectLibre-Export) -> nach anhaenge/ hochladen,
%%              Variante B auskommentieren und die \anhangpdf-Zeile aktivieren.
%%  Variante B: Gantt-Diagramm direkt in LaTeX (pgfgantt) – Gerüst unten anpassen.
%%%%%%%%%%%%%%%%%%%%%%%%%%%%%%%%%%%%%%%%%%%%%%%%%%%%%%%%%%%%%%%%%%%%%%%%%%%%%%%%

%% ---- Variante A ----
%\anhangpdf{Arbeitsplan (Einlage 1)}{anh:arbeitsplan}{anhaenge/arbeitsplan.pdf}

%% ---- Variante B ----
\chapter{Arbeitsplan (Einlage 1)}
\label{anh:arbeitsplan}

\hinweis{Gerüst mit den festen Terminen des Projektpraktikums. Arbeitspakete, Zeiträume und Zuständigkeiten
(\IhrName{} / \Partner{} / gemeinsam) an euer Projekt anpassen; feiner untergliedern; Parallelität sichtbar
machen. Ein Tag = 0{,}58\,mm. Balken: \texttt{\textbackslash ganttbar\{Name\}\{Start\}\{Ende\}}, Termine:
\texttt{\textbackslash ganttmilestone\{Name\}\{Datum\}}.}

\begin{figure}[htbp]
\centering
\begin{ganttchart}[
    time slot format=isodate,
    x unit=0.58mm, y unit chart=0.5cm, y unit title=0.55cm,
    hgrid, vgrid={*6{draw=none},dotted},
    title label font=\scriptsize, bar label font=\scriptsize,
    group label font=\scriptsize\bfseries, milestone label font=\scriptsize\itshape,
    bar height=0.55, group height=0.35,
    bar/.append style={fill=blue!25},
    group/.append style={fill=black!60},
    milestone/.append style={fill=red!70},
  ]{2026-10-12}{2027-03-21}
  \gantttitlecalendar{year, month=shortname} \\
  \ganttgroup{Planung}{2026-10-12}{2026-10-30} \\
  \ganttbar{Themen- und Schwerpunktwahl}{2026-10-12}{2026-10-15} \\
  \ganttbar{Motivationsformular}{2026-10-15}{2026-10-21} \\
  \ganttbar{Arbeitsplan erstellen}{2026-10-15}{2026-10-28} \\
  \ganttgroup{Literatur}{2026-10-19}{2026-11-11} \\
  \ganttbar{Literaturrecherche}{2026-10-19}{2026-11-04} \\
  \ganttbar{Paper auswerten (Einlage 2)}{2026-10-28}{2026-11-11} \\
  \ganttgroup{Experiment}{2026-10-22}{2027-01-15} \\
  \ganttbar{Aufbau planen, Material}{2026-10-22}{2026-11-12} \\
  \ganttbar{Aufbau und Test}{2026-11-12}{2026-12-03} \\
  \ganttbar{Messungen}{2026-12-03}{2027-01-08} \\
  \ganttbar{Puffer (Weihnachten)}{2026-12-21}{2027-01-06} \\
  \ganttgroup{Auswertung}{2026-12-10}{2027-01-20} \\
  \ganttbar{Auswertung, Fehlerrechnung}{2026-12-10}{2027-01-15} \\
  \ganttbar{Messergebnis (Einlage 4)}{2027-01-08}{2027-01-20} \\
  \ganttgroup{Vortrag}{2027-01-08}{2027-01-28} \\
  \ganttbar{Folien und Probevortrag}{2027-01-08}{2027-01-27} \\
  \ganttgroup{Portfolio}{2026-10-15}{2027-03-15} \\
  \ganttbar{Laborjournal (laufend)}{2026-10-15}{2027-01-28} \\
  \ganttbar{Vortrags-Einlagen 5 und 6}{2027-01-29}{2027-02-21} \\
  \ganttbar{Einleitung, Gesamtreflexion}{2027-02-15}{2027-03-08} \\
  \ganttbar{Korrektur, Endabgabe}{2027-03-01}{2027-03-15} \\
  \ganttmilestone{Abgabe Arbeitsplan}{2026-10-28} \\
  \ganttmilestone{Abgabe wiss.\ Arbeit}{2026-11-11} \\
  \ganttmilestone{Abgabe Protokoll}{2026-12-09} \\
  \ganttmilestone{Abgabe Messergebnis}{2027-01-20} \\
  \ganttmilestone{Vortrag (Prüfung)}{2027-01-28} \\
  \ganttmilestone{Endabgabe Portfolio}{2027-03-15}
\end{ganttchart}
\caption{Arbeitsplan von \IhrName{} und \Partner{} für das Projekt \enquote{\Projekt}
(gemeinsam erstellt). Rote Rauten: Abgabetermine.}
\label{fig:arbeitsplan}
\end{figure}


\end{document}
